\chapter{Conclusion}\label{ch:conclusion}

\section{Summary \& Findings}

This thesis addressed the reconstruction of scenario-specific Interaction
Models from \Cpp\ codebases and execution traces in event-driven systems.
Static analysis alone tends to produce over-approximate, scenario-agnostic
graphs; trace-based analysis alone tends to under-complete them, since
runtime instrumentation is rarely exhaustive. No prior approach appeared
to address publish-subscribe topology recovery and scenario-scoped
call-graph querying together, or to target \sysml\ as the output
formalism.

The approach (\cref{ch:methods}) combines a static Neo4j call graph
with a scenario-scoped runtime trace, filters the trace through a
Neo4j-based reachability filter (condition~C5), and projects the
result into \sysml\ interaction diagrams.

The approach was developed and evaluated on \emph{CPSCore}, an
open-source event-driven \Cpp\ framework, across four found-in-the-wild
scenarios drawn from its existing test suite, including two instances of
the same underlying scenario (its happy path and its timeout branch).

The three hypotheses tested to answer the research questions were
confirmed, with an important qualification on \emph{when} and
\emph{how well} each effect holds:

\begin{description}
  \item[H1 (confirmed).] \textbf{C5 (guided~+~reachability) and C6
    (full-trace~+~reachability) achieve perfect F1 on both interaction-set
    and execution-order metrics across all four scenarios}
    (macro-F1~$=$~1.000\,/\,1.000). The reachability filter is the key
    mechanism: C5 and C6 differ only in instrumentation strategy yet
    score identically, confirming the filter is sufficient regardless of
    how the trace is collected. \textbf{C4 (guided-only, 0.914\,/\,0.846) is
    the strongest baseline without the Neo4j-based reachability filter}, substantially outperforming
    unfiltered tracing (C3, 0.530\,/\,0.428); the gap is explained entirely
    by intra-component false positives in C3's input.
    \textbf{C2 (static-only) is structurally capped on S3} by indirect
    thread dispatch (macro-F1~$=$~0.756): no static call site exists for
    the \\texttt{Utilities}$\\to$\\texttt{Logging} edge, so the trace-based
    conditions are the only ones that recover it.
    \textbf{LLM-augmented conditions (C7--C9) performed worse than
    their non-augmented counterparts}: C7 recovers most of C2's recall but
    still fails on S3; C9 degrades C4's already-precise output
    (0.783~vs.~0.914), confirming that LLM post-processing is
    counterproductive when the prior is clean.

  \item[H2 (confirmed).] Code-graph-guided instrumentation eliminates
    all measured noise (50\% fewer retained nodes, 70\% fewer edges)
    without losing coverage, and a blinded six-rater panel rated the
    guided output strictly higher in 47 of 54 comparisons, with the one
    reversal traceable to a specific trade-off: collapsing repeated
    property-read calls into a single edge.

  \item[H3 (supported).] Diagram-inclusive conditions were never worse
    than their non-diagram counterparts on the scenario-comprehension
    task tested, and clearly outperformed source code alone (the
    weakest condition, missing a component and interactions that trace
    or diagram evidence made explicit). The reconstructed diagram alone
    tied with runtime traces alone on interaction recall (0.778),
    sharing the same attribution error on one scenario; only combining
    source code with runtime traces reached perfect recall, and adding
    a diagram on top raised completeness further without cost to
    recall. A reconstructed diagram is, on this task, a useful
    complement to the underlying evidence rather than a strictly
    sufficient replacement for it on its own.
\end{description}

As an exploratory complement, a hands-on session and interview study with
four engineers on a second, much larger industrial codebase
(\cref{sec:practitioner}) offers encouraging evidence that the static
half of the approach holds at industrial scale, though the session
exercised only static extraction and graph querying, with no runtime
tracing involved. The key scenario there was a proxy-leak defect whose
root cause was an implicit event-driven interaction, invisible without
reconstructing the call chain, resolved using only the static graph,
consistent with static extraction remaining useful for fault-path
investigation independent of runtime access. The industrial system's
architecture (proxy-based, virtual interfaces, dynamic service binding)
is the kind where static analysis is more likely to carry its own
independent errors, which H1 would predict gives the static-graph
reachability filter more trace events to suppress
(\cref{sec:cpscore-scale,sec:h1-c4c2}); since no dynamic
evidence was collected there, this remains a motivated prediction rather
than a measured result.

Taken together, a single pattern runs through the H1 results
(\cref{sec:discussion}): additional evidence helps when it disagrees
productively with the primary source, by independently confirming which
of its edges are genuine, and is otherwise redundant -- and even a
successfully combined, high-precision reconstruction is not automatically
a sufficient basis for a downstream comprehension task, which H3 shows
needs the underlying source and trace evidence directly. This generates a
falsifiable prediction for future work: on systems with cross-module
virtual dispatch, static analysis should carry independent errors even
on happy-path scenarios, giving intersection more independent error to
cancel across a wider range of scenarios than the single instance
(\texttt{std::function} dispatch) exercised here
(\cref{sec:future-work}).

\section{Practical Implications}

The pipeline lets an engineer working on an event-driven \Cpp\ system
obtain a scenario-specific \sysml\ interaction diagram through
natural-language queries against the property graph, rather than
manually tracing call chains, cross-referencing headers and correlating
log output over several working days. \Cref{sec:utility} quantifies this
on CPSCore (2,574 interactions extracted in a single query, nine caller
functions instrumented automatically, without any manual call-graph
reading); switching to a new target requires only re-running the static
extractor, since the skills query the live schema rather than hardcoded
labels. The practitioner session (\cref{sec:practitioner}) suggests this
transfer cost is low in practice: four engineers on a much larger
industrial codebase used the tool productively within a single session,
and rated the structural and dependency diagrams consistently useful for
architectural comprehension.

\section{Contributions}

\paragraph{Research contributions.}
\begin{enumerate}
  \item \textbf{A static property graph and a separate trace pipeline}
    whose interaction edge sets are compared while keeping track of
    which evidence each interaction came from: static, dynamic, or
    both. Keeping the sources separate but comparable is what lets the
    evaluation report exactly where each edge originated
    (\cref{sec:pg-schema}).

  \item \textbf{A scenario-scoping method} that trims the call graph to
    the interactions relevant to one scenario, using plain
    scenario-scoped Cypher queries, without re-running the analysis from
    scratch (\cref{sec:combine-scope}).

  \item \textbf{An evidence-based account of when combining static and
    dynamic analysis helps}, and when it does not, summarised above for
    H1--H3 (\cref{sec:h1-results,sec:h2-results,sec:h3-holistic}).
\end{enumerate}

\paragraph{Engineering artefacts.}
\begin{enumerate}\setcounter{enumi}{3}
  \item A \textbf{small, reversible instrumentation tool}
    (\texttt{csp\_matcher}) that inserts and removes trace probes
    automatically, driven by a simple find/replace specification
    (\cref{sec:csp-matcher}).

  \item A \textbf{set of twelve natural-language agent skills} that
    together run the whole pipeline, from source code to \sysml\ diagram,
    on any codebase without code changes, every one following a fixed
    rule set given its inputs
    (\cref{sec:skill-design,sec:evidence-synthesis}).

  \item A \textbf{way to recover publish-subscribe wiring} that ordinary
    call-graph analysis cannot see, by matching subscription and
    publication call sites that share a channel identifier
    (\cref{sec:macro-resolution}).
\end{enumerate}

\section{Limitations \& Threats to Validity}

The most significant limitation is scale: all three hypotheses were
evaluated on a single open-source system ($\approx$13\,kLOC) across four
scenarios. CPSCore's absence of cross-module virtual dispatch means its
static call graph disagrees with the trace on \emph{which} edges are
genuine in only one scenario (an indirect thread dispatch), which is
presumably why intersection's recall shortfall is concentrated there
rather than spread across every scenario; in a larger system with
virtual dispatch or dynamic plugin loading, where static analysis is
more likely to carry its own independent errors on a wider range of
scenarios, intersection's precision advantage would have more
disagreement to cancel, but its recall shortfall could also be more
widespread. The practitioner session on an industrial-scale codebase
offers encouraging qualitative evidence of transfer, but does not
substitute for a controlled replication.

H3's diagram-inclusive conditions were shown diagrams rendered from the
guided-only condition~(C4) rather than from the
guided+reachability condition~(C5) (\cref{sec:h3-holistic}). Its interaction-recall figures also
depend on a manual correction to the automated scorer, which cannot
tell a genuinely-attributed interaction from one whose components merely
co-occur in the response text. The H2 six-rater panel preserved effect direction on
every dimension, though a larger sample would strengthen confidence in
magnitude. Ground-truth diagrams are inherently subjective, and the
edge-set metric does not capture message content or count even where
the execution-order metric captures ordering.

\section{Future Work}\label{sec:future-work}

\begin{description}
  \item[Controlled evaluation on a second system.]
    The practitioner session (\cref{sec:practitioner}) suggests the
    static approach transfers to larger, industrial codebases, but only
    qualitatively. A natural first step would be applying the full,
    pre-registered H1/H2/H3 tests to that codebase or to other
    event-driven \Cpp\ systems (\eg using ROS~2 or MQTT for
    publish-subscribe) to turn this encouraging signal into a controlled
    generalisability result, ideally on a system with cross-module
    virtual dispatch so that static analysis carries independent errors
    on more than the single scenario type observed here.
 
  \item[Diagram labelling and interactivity.]
    The practitioner interviews (\cref{sec:interview-findings})
    independently converged on two requests: richer, less abstract
    event/message labels on sequence diagrams (\eg inline source class
    and line number on hover), and interactive exploration (search,
    filter, collapse/expand) rather than static, one-shot diagrams. Both
    are natural extensions of the source-location data the pipeline
    already carries (\cref{sec:utility}).

  \item[Temporal ordering accuracy.]
    The current evaluation abstracts over message ordering. A metric
    accounting for the relative order of interactions would give a more
    complete picture.

  \item[Incremental updates.]
    As the codebase evolves, the \ac{PG} should ideally be updatable
    incrementally rather than rebuilt from scratch. Neo4j's support for
    live graph updates makes this tractable, though the
    scenario-scoping invalidation logic would need careful design.

  \item[LLM-assisted diagram naming.]
    Function names in the generated diagrams are currently taken
    verbatim from source. A post-processing step using an \ac{LLM} could
    produce more readable message labels.

  \item[Integration with Renaissance DB.]
    The \ac{GDI} project's architecture model database (Renaissance DB)
    could serve as a richer source of module-level metadata, potentially
    reducing the number of unresolved black-box components.
\end{description}

\section{Reproducibility}

The complete source code, experiment data, skill definitions, and scoring
scripts used in this thesis are available at:

\begin{center}
\url{https://github.com/prathikanand7/masters-thesis-work.git}
\url{https://ci.tno.nl/gitlab/anand_krishnan.prathik-tno/cpscore.git}
\end{center}
